\subsection{Exchanging theta integrals of complementary ranks}\label{sec:exchange-proof}

The determinant comparison uses an identity between integrals along two different
groups, viewed as functions on the same group. Write \(n,d_+,d_-\) for the ranks of
\(W,V_+,V_-\), with \(n=d_++d_-\) and \(d_+-d_-=2e\ge0\).
The number \(q\) of negative signs of \(W\) is the sum of the numbers for \(V_+\)
and \(V_-\). If \(\mu_\tau\) is the weight of \(\nu\) at an infinite place,
the range is \(2q+2\le n\le2\mu_\tau+2q\).
The character \(\nu\) is tamely ramified at a nonempty finite set \(S\) of split
places disjoint from its conjugate. Write \(\chi\) for its ramified part and
\(L^S(s,\nu)\) for the \(L\)-function of its base change with the factors
at \(S\cup\bar S\) omitted. Put
\begin{equation}\label{eq:exchange-lstar}
 L_*(e)=\left(\prod_{w\in S}\mathrm Nw\right)^{e^2}
             \prod_{j=1}^{e}|L^S(j,\nu)|^2,\qquad
 L_*(-e)=L_*(e)^{-1}.
\end{equation}
For an integrally invertible matrix \(g\), let \(R(g)\) consist of the last
\(d_+\) rows of \(g^{-1}\) and \(C(g)\) of the first \(d_-\) columns of \(g\).
The finite data \(\Phi,\Psi\) are supported on these respective matrices
and satisfy \(\Phi(R(g))=\chi(\det g)\Psi(C(g))\).

Given the real datum \(y\) on the \(V_-\) side, one can choose the datum \(x\)
on the \(V_+\) side before choosing the modulus and the character.
Writing the two integrals as \(I_+(x,\Phi;g)\) and \(I_-(y,\Psi;g)\), we have
\begin{equation}\label{eq:exchanged-integrals}
 I_+(x,\Phi;g)=c\,I_-(y,\Psi;g),\qquad
 |c|^2=C L_*(e)>0.
\end{equation}
Here \(g\in U(W)(\mathbb A)\) and \(C>0\) is independent of the modulus and character.
In the Weil representation the parameter for \(W\) is \(\gamma=\nu^{-1}\upsilon^{n/2}\).
The two parameters \(\alpha,\alpha'\) for the partner spaces satisfy
\(\alpha\alpha'=\gamma\). This is equivalent to the product of the
rank-dependent splitting characters being the inverse of the base change of \(\nu\).
\leanlinks{\leansrc{New/Determinant/Exchange/TheoremForPartners}{Hypotheses and normalization of the exchange identity}}

\paragraph{Comparison after adjoining hyperbolic planes}
Take adjoints and reverse the signs of \(V_-\) to write both kernels as kernels
from the central space. Put \(p=d_-\), \(W_1=W\perp H\), and \(X=W_1\perp H^p\).
The large partner is \(V_-^{\mathrm{opp}}\perp W_1\) and the small partner
is \(V_-^{\mathrm{opp}}\). First extend \(y\) to a Schwartz vector \(y_1\)
whose zero slice over \(W\) is \(y\). Choose a Schwartz vector \(v_1\)
for \((W_1,W_1)\) as a second input. At the ramified places use the standard
data. For an integral auxiliary matrix \(z\), the row datum is \(\chi^{-1}\)
of the determinant of the square matrix obtained by stacking \(z\) and the
variable. The column datum uses the determinant of their product. Both are
zero when the determinant is a nonunit.

In the mixed model for \(X\), take slices along the \(p\) isotropic columns.
At an infinite place write \(x_0=\binom{I}{0}\) for the standard frame and \(J\)
for the matrix of the Hermitian form of the large partner. Near \(x_0\), the
span of \(x_0\) can be moved to that of a frame \(x\) by a smoothly varying
unitary matrix. Indeed, let \(w=(x_0^*Jx_0)^{-1}x_0^*\) and put
\[
 A(B)=(Bx_0w-x_0wB^*)J,\qquad
 \mathcal T(B)=\exp(A(B))(B-A(B)).
\]
Then \(A(B)^*J+JA(B)=0\) and \((B-A(B))x_0=x_0wB^*Jx_0\).
As \(A(I)=0\), the derivative of \(\mathcal T\) at \(I\) is the identity.
Apply its local inverse to \(I+(x-x_0)wJ\) to obtain \(B\).
Then \(k(x)=\exp(A(B))\) carries the span of \(x_0\) onto the span of \(x\).

Choose a nonnegative smooth cutoff \(\sigma\) supported in this neighbourhood
and equal to \(1\) at \(x_0\). Use the projection of \(x_0\) onto the span of
\(x\), namely \(z(x)=x(x^*Jx)^{-1}x^*Jx_0\).
The inverse is needed only in the cutoff neighbourhood and can be smoothly
extended away from it. On the large side act on \(y_1\otimes v_1\) by
\(k(z(x))\). Multiply by \(\sigma(z(x))\), by the cutoff at the moved frame
\(k(z(x))^{-1}x\), and by the \(2\mu_\tau\)-th power of the conjugate
determinant of its upper square block at each infinite place.
On the small side use the cutoff of the standard frame and
\(|\det B_\tau|^{2(j+\mu_\tau)}\) in the square variable \(B_\tau\), where
\(2j=\operatorname{rk}W_1\). Normalize this slice so that the zero slice
in the opposite polarization is \(y_1\). The zero slice is obtained by taking
the Fourier transform at zero in the frame variable, and the integral of the
cutoff with the determinant powers is positive. Dividing by that integral
gives the required normalization. Compact support in the frame variables
and the Schwartz property in the orthogonal variables give the required
Schwartz vectors.

For \(h\in U(V_-^{\mathrm{opp}}\perp W_1)(E_\infty)\) with nonzero cutoff at
\(z(h^{-1}x_0)\), the matrix \(h k(z(h^{-1}x_0))\) preserves the span of \(x_0\).
It is therefore block diagonal. Absorb its two blocks into the variables
of integration on \(U(V_-^{\mathrm{opp}})\) and \(U(W_1)\).
The moved frame is now \(\binom{B_\tau}{0}\), and the theta functional
factors. Integration over \(U(W_1)\) contributes
\[
 I_0=\int_{[U(W_1)]}
    \theta_{W_1,W_1}\bigl(v_1,\chi^{-1}\circ\det\bigr)(1,h_1)\,dh_1 .
\]
The other factor is the small-side slice. The ratio of the determinant
factors is
\[
 \frac{1}{\nu(b_\infty/\bar b_\infty)|b_\infty|_{\mathbb A_E}^{j}},
\]
where \(b_\infty\) has infinite part \(\det B_\tau\) and finite part \(1\).
This cancels the inverse factor in the Levi action on the kernel.
\leanlinks{\leansrc{New/Determinant/Exchange/RealVectors}{Moving the frame smoothly and computing the slice period}}

Taking a central Fourier coefficient restricts the theta sum to the orbit
of frames with the specified Gram matrix. A real signature absent on the
small side contributes zero on the large side by the support of the cutoff.
If the finite lattice is self-dual and the signatures agree, the local-global
principle for Hermitian forms and Witt's theorem reduce the index to the
standard Gram matrix. The stabilizer of that frame on the large side is
\(U(W_1)\) and that on the small side is trivial. Unfold the large-side
integral in stages through \(U(V_-^{\mathrm{opp}})\times U(W_1)\) and use the
slice calculation. The ratio of the two coefficients is the product of
\(I_0\), a measure constant, and character values of modulus one.
The measure constant is nonzero because the nonnegative cutoff is positive
near the identity.

The self-dual lattice tests suffice as follows. At a ramified place, translate
the column datum by the root group of the additional hyperbolic line, with
matrix \(A\) ranging modulo \(\varpi_w\). On the central Fourier coefficient
of the difference this gives the multiplier
\[
 \sum_{A}\chi_w^{-1}(\det A)\,\psi_w(\operatorname{tr}(TA)).
\]
For a primitive character the matrix Gauss sum vanishes at singular \(T\).
At invertible \(T\), the change of variable \(A\mapsto TA\) gives the sum
at the identity times a determinant character. Thus only indices invertible
at the ramified places have to be tested. At the other finite places, the
shears through the adjacent hyperbolic plane move the partial means of the
central coefficient. They fix the block reading the left inverse in the
standard column datum. Strong approximation moves the coefficient to a
point whose Gram matrix and inverse are integral, where there is an integral
orthonormal basis. The self-dual comparison applies there. Central Fourier
inversion and parabolic decomposition give proportionality on the standard data.
\leanlinks{\leansrc{New/Determinant/Exchange/LargeTriple}{The Fourier coefficients on Gram orbits}
  \leanlinksep \leansrc{New/Determinant/Exchange/StableSpacesStandardData}{Matrix Gauss sums and self-dual lattice tests}}

\paragraph{Rank descent by constant terms}
Remove corresponding hyperbolic planes one at a time. Let \(u\) be the rank
of the remaining large partner minus half the remaining central rank.
Unfolding the large-side constant term in the mixed model gives Tate's
integral along the isotropic line at \(s=u+1\). The determinant character of
the standard datum kills the zero term and the Fourier-transformed zero
term. At the ramified places the integral gives the Gauss sum of the unit
character, and at the unramified places the vacuum gives a geometric series.
At infinity one integrates the slice on the isotropic cone over the maximal
compact group and the infinite ideles. For \(u\ge0\) the constant term
is \(\gamma_u\) times the lower theta integral with the cone datum at \(s=u+1\).
For \(u<0\) the other side of Poisson summation gives the dual cone datum
at \(s=-u\). In both cases, with \(a_u>0\) independent of the character,
\begin{equation}\label{eq:exchange-descent}
 |\gamma_u|^2 L_*(u)=a_u L_*(u+1).
\end{equation}
For positive \(s\), the Tate factor has modulus, up to a fixed factor,
\(|L^S(s,\nu)|\prod_{w\in S}(\mathrm Nw)^{s-1/2}\).
The difference \((u+1)^2-u^2=2(u+1)-1\) gives the formula for \(u\ge0\);
the inverse convention for \(L_*\) gives it for \(u<0\).
On the column side, primitivity of the full column block kills the
contribution of a nonzero isotropic column. The constant term is the
integral of the zero slice times a fixed vacuum factor \(\kappa\ne0\).
Thus removing one plane changes the proportionality factor \(c\) to
\(c\kappa/\gamma_u\). Equation~\eqref{eq:exchange-descent} preserves the
ratio to \(L_*\) up to a constant independent of the character.
Removing all \(p+1\) added planes restores the required column datum.
The squared modulus of the resulting factor is \(L_*(e)|I_0|^2/L_*(j)\)
times a fixed positive constant.
\leanlinks{\leansrc{Lib/SiegelWeil/Tower/LargeSideConstantTerm}{The Tate integral for a constant term and its local factors}
  \leanlinksep \leansrc{New/Determinant/Exchange/RemovingPlanes}{Descent by cone integrals and zero slices}}

\paragraph{The equal-rank normalization and nonvanishing}
The equal-rank integral with finite datum \(\chi^{-1}\circ\det\) is constant.
At the ramified places the general linear group fixes the datum because
the column block is empty. Elsewhere the standard level fixes it. Strong
approximation, together with the constant term calculation along the
isotropic line, shows that the dependence on the line variable of the
Levi also disappears. Let \(R_{W_1}(v_1)\) denote the squared modulus of this
integral divided by \(L_*(\operatorname{rk}W_1/2)\).

First take \(J_0\) positive definite of rank \(2k_0\) and \(W_1=H\perp J_0\).
If \(R_{W_1}\) were zero for every datum at a fixed character, the preceding
Fourier coefficient argument would give vanishing on all standard data
on the large side for \((H^{2k_0}\perp W_1,J_0^{\mathrm{opp}}\perp W_1)\).
The vanishing descends one plane at a time. For any datum \(v\) below,
place a smooth cutoff on a small region of the isotropic cone and take
the slices there to be compact-group translates of \(v\). The Tate integral
above separates a strictly positive Mellin factor \(B(s)\) and the lower
integral \(I(v)\). The product formula for \(s\ge1\) extends analytically
to \(s>-k_0\). With weights \(\mu_\tau\ge k_0\), the functional equation
has no vanishing infinite-place factor at the relevant integers in this range.
The Euler product for \(s>1\), nonvanishing at \(s=1\), and the functional
equation give nonvanishing of the \(L\)-factor at every integer \(s>-k_0\).
Thus evaluating at \(u+1>-k_0\) makes \(I(v)=0\) whenever the upper
integral vanishes. After \(2k_0+1\) steps one reaches \((J_0,0)\), where the
vacuum integral is the strictly positive volume of the rank-zero quotient.
This contradicts vanishing.

Now apply the exchange successively to \((J_0,J_0,0)\) and \((J_0,0,J_0)\).
The squared moduli of the resulting factors \(c_1,c_2\) are proportional
to \(L_*(k_0)R_{W_1}(v_1)\) and \(L_*(k_0)^{-1}R_{W_1}(v_1)\), respectively,
with constants independent of the character. Integrate the intermediate
\((J_0,J_0)\) kernel over both groups and exchange the order by Fubini.
This identifies \(c_1c_2\) with the value of the real datum in rank zero,
independent of the character. The \(L_*(k_0)\) factors cancel, so
\(R_{W_1}(v_1)^2\) is independent of the character. Choose \(v_1\) by
nonvanishing for one character. Then \(R_{W_1}(v_1)>0\) has the same value
for all characters.
For general signs, add hyperbolic planes. At each stage the upper equal-rank
integral equals its constant term and the lower datum can be realized by
the cone integral. By \eqref{eq:exchange-descent} this multiplies the ratio
by a fixed positive constant.
\leanlinks{\leansrc{New/Determinant/Exchange/ConeImage}{Realizing prescribed data by cone integrals}
  \leanlinksep \leansrc{New/Determinant/Exchange/VanishingDescends}{Descending vanishing by constant terms}
  \leanlinksep \leansrc{New/Determinant/Exchange/TheoremForPartners}{Nonvanishing and character independence of the equal-rank ratio}}

\paragraph{Passage to all matched data}
The difference of the two integrals is an equivariant linear map for
left translations at the ramified places. Its kernel contains the
standard column datum \(P(Y)=\chi^{-1}(\det zY)\).
Choose integral \(x_0\) with \(z x_0=I\), and choose \(a\) so that \(\Psi\)
is constant on residue classes modulo \(\varpi^a\). The matrix
\(d=x_0z+\varpi^{-a}(I-x_0z)\) fixes the span of \(x_0\) and expands \(\ker z\).
When \(P(dgY)\ne0\) we have
\(Y\equiv g^{-1}x_0(z\,dgY)\pmod{\varpi^a}\).
The determinant character for the right action gives
\(\Psi(g^{-1}x_0)P(dgY)=|P(dgY)|\Psi(Y)\).
Average over \(g\in\mathrm{GL}_n(\mathcal O_S)\). The left side belongs
to the kernel. The mean of \(|P(dgY)|\) is constant on the support of
\(\Psi\) and positive near \(Y=x_0,g=1\). Hence \(\Psi\) itself belongs to
the kernel. Apply this argument to data modulo \(\mathfrak m\) by lifting
them to the local integer rings. Undoing the adjoint on the small side
transposes the data and scales by an integral unit of norm \(-1\), which
contributes a scalar of modulus one. The relation \(\alpha\alpha'=\gamma\)
then inserts the same factor \(\alpha(\det g)^{-1}\) on both sides.
This proves \eqref{eq:exchanged-integrals}.
\leanlinks{\leansrc{New/Determinant/Exchange/StableSpacesStandardData}{Generating the matched data by averaging translates}
  \leanlinksep \leansrc{New/Determinant/Exchange/TheoremForPartners}{Adjoints and changing the splitting characters}}

\paragraph{Composing two lifts}
Take \(W=U_1\perp U_2\), \(V_+=A\), and \(V_-=A'\).
Use the two column blocks of row datum \(x\) in the respective kernels,
and write \(\Theta_x\) for the lift \(U_1\to A\) followed by the transposed
lift \(A\to U_2\). Write \(\Theta'_y\) for the composition through \(A'\)
using the two row blocks of column datum \(y\).
Fix forms \(f_1,f_2\) on \(U_1,U_2\) and pair the resulting lift of \(f_1\)
with \(f_2\). For the edge characters \((a_i,b_i)\) and \((a_i',b_i')\)
the product condition gives \(a_1a_2=a_1'a_2'=\gamma\) and
\(\nu_0=b_1/b_2=b_2'/b_1'\). Multiply the exchange at \((g_1,g_2)\) by
\(f_1(g_1)\overline{f_2(g_2)}\nu_0(\det g_2)\).
The last character cancels the difference in characters when the kernel
on each path is split into the product of its two edge kernels.

On the \(A'\) side the real datum starts as the product of the two
Fock polynomials. The point data also split by the two row blocks.
Fubini and transposition identify the integral with
\(\langle f_2,\Theta'_y f_1\rangle\). On the \(A\) side the real datum
is the vector chosen by exchange. First decompose a polynomial into
a sum of products in the two column blocks. For each product, real
reproduction on each edge writes its lift as a real convolution of the
prescribed edge lift. In the pairing the two convolutions give a scalar
depending only on the real data. The two forms are rapidly decreasing
for the types of the edges, so the integral is continuous in the moment
norms of the Schwartz data. The proportionality to all pairings is a
closed condition, so the integral for the chosen real vector is
\(\kappa_0\langle f_2,\Theta_x f_1\rangle\).
Summing both sides with the finite weights and testing against every
\(f_2\) gives
\begin{equation}\label{eq:exchange-square}
 \sum_y\Psi(y)\Theta'_y=\frac{\kappa_0}{c}\sum_x\Phi(x)\Theta_x .
\end{equation}
This uses \(c\ne0\). The scalar \(\kappa_0\) is independent of the twist,
and the identity at this stage allows \(\kappa_0=0\).
\leanlinks{\leansrc{New/Determinant/Exchange/OneSquare}{Exchanging the two paths with transposed lifts}}
